\documentclass[10pt]{article}
\usepackage[margin=0.76in]{geometry}
\usepackage[T1]{fontenc}
\usepackage{lmodern}
\usepackage{amsmath,amssymb,amsthm}
\usepackage[hidelinks]{hyperref}
\pdftrailerid{C115McMillanRational20260824}
\pdfinfoomitdate=1
\pdfsuppressptexinfo=15
\newtheorem{proposition}{Proposition}
\title{A Reversible Rational McMillan Map:\\Exact Pole Exclusion and Low-Period Monodromy}
\author{Anonymous Route-A report}
\date{}
\begin{document}
\maketitle
\begin{abstract}
We examine the rational recurrence
$M(x,y)=(-4x/(1+x^2)-y,x)$.  Exact symbolic calculation verifies its rational
inverse, coordinate-swap reversor, unit Jacobian determinant, and quartic first
integral.  The valid fixed locus over $\mathbb C$ contains one real and two
nonreal points, while $(1,-1)\leftrightarrow(-1,1)$ is a genuine real
primitive two-cycle.  We expose and remove the factor $(x^2+1)^2$ introduced
by denominator clearing: its roots are poles, not periodic points.  The
cycle's two-step monodromy is $-I_2$.  These are finite low-period
certificates, not a global orbit classification or transfer determinant.
\end{abstract}

\section{Rational map and invariant}
We freeze the McMillan/QRT-type family
\[
 M_\mu(x,y)=\left(\frac{2\mu x}{1+x^2}-y,x\right),\qquad \mu=-2.
\]
Over $\mathbb C$ the forward map excludes $x^2+1=0$.  Its rational inverse is
\[
 M^{-1}(x,y)=\left(y,-\frac{4y}{1+y^2}-x\right).
\]
For the involution $S(x,y)=(y,x)$, direct composition gives
$SMS=M^{-1}$.  Differentiation gives
\[
 DM(x,y)=\begin{pmatrix}
 4(x^2-1)/(1+x^2)^2&-1\\1&0
 \end{pmatrix},\qquad \det DM=1.
\]
Substitution and cancellation also give
\[
 I\circ M=I,\qquad I(x,y)=x^2y^2+x^2+y^2+4xy.
\]

\section{Fixed points and a two-cycle}
The fixed equations force $y=x$ and reduce to
\[
 -\frac{2x(x^2+3)}{x^2+1}=0.
\]
Thus the valid fixed points are $(0,0)$ and
$(\pm i\sqrt3,\pm i\sqrt3)$, with matched signs.  Only the origin is real.
Direct substitution gives the distinct real cycle
\[
 q_+=(1,-1)\longmapsto q_-=(-1,1)\longmapsto q_+,
 \qquad I(q_\pm)=-1.
\]
Both cycle denominators equal $2$, so neither point is a pole.

For the second iterate, eliminating $y$ from the two cleared numerators gives
\[
 -8x(x-1)(x+1)(x^2+1)^2(x^2+3).
\]
The roots of $x^2+1$ cannot be tested as orbit points because the first
application of $M$ is already undefined.  Removing this pole factor leaves
$x(x-1)(x+1)(x^2+3)$: the factors $x(x^2+3)$ recover the fixed locus and
$x^2-1$ gives the primitive cycle.
\begin{center}
\small
\begin{tabular}{llll}
candidate & type & $1+x^2$ & status\\ \hline
$0$ & real fixed & $1$ & valid\\
$\pm i\sqrt3$ & nonreal fixed & $-2$ & valid\\
$\pm1$ & real two-cycle & $2$ & valid\\
$\pm i$ & cleared-denominator root & $0$ & excluded\\
\end{tabular}
\end{center}

\section{Local period-two monodromy}
At either cycle point the first derivative is
\[
 J_\pm=\begin{pmatrix}0&-1\\1&0\end{pmatrix}.
\]
In chronological order the derivative of $M^2$ is therefore
\[
 P_2=J_-J_+=-I_2,\qquad
 \det P_2=1,\quad \operatorname{tr}P_2=-2.
\]
Consequently
\[
 \det(\lambda I-P_2)=(\lambda+1)^2,
 \qquad \det(I-zP_2)=(1+z)^2.
\]
As a one-step control, the real fixed point has
\[
 J_0=DM(0,0)=\begin{pmatrix}-4&-1\\1&0\end{pmatrix},\qquad
 \det(I-zJ_0)=z^2+4z+1.
\]
The unequal control polynomials prevent a one-step fixed-point linearization
from being conflated with the chronological two-step cycle product.

\begin{proposition}
The inverse, reversor, invariant, fixed points, pole exclusions, primitive
cycle, and displayed monodromy identities are exact on their stated domains.
\end{proposition}
\begin{proof}
An independent checker recomputes both rational compositions, differentiates
the map, substitutes into $I$, recomputes the raw resultant, tests every
candidate denominator and orbit arrow, and multiplies the Jacobians without
importing the evidence producer.
\end{proof}

\section{Route-A boundary}
The two displayed $z$-polynomials are characteristic data of finite local
derivatives.  They have neither orbit weights nor a trace identity over a
global state space.  In particular, $(1+z)^2$ belongs only to the local
derivative of $M^2$ along one cycle.  No transfer operator, finite transfer
owner, function space, or tail bound has been constructed.  The result is
\texttt{A1\_PARTIAL\_CERTIFIED}, \texttt{A2\_FAIL},
\texttt{A3\_NOT\_ADDRESSED}, and \texttt{A4\_FAIL}.  We do not claim a
complete orbit atlas, global integrability classification, analytic Fredholm
determinant, arithmetic data, or Route-B object.  The scope firewall is
\texttt{NO\_BAD\_EULER\_OR\_ROOT\_NUMBER}.
\end{document}
